\notechapter{N-20221002}{Finite Fields and Polynomial Factorization}

\section{Why finite fields matter}

This note is a more systematic essay on finite fields and factorization.  The central question is: how do polynomials factor when the coefficients live in a finite field rather than in \(\mathbb Q\), \(\mathbb R\), or \(\mathbb C\)?

This question is not only theoretical.  It appears in coding theory, cryptography, computational algebra, and the construction of algebraic extensions.  The note correctly starts from the definition of a field, because factorization depends heavily on the coefficient domain.

\section{Fields}

A field \(F\) is a set with addition and multiplication such that:
\begin{enumerate}[(i)]
  \item \((F,+)\) is an abelian group with identity \(0\);
  \item \(F\setminus\{0\}\) is an abelian group under multiplication with identity \(1\);
  \item multiplication distributes over addition.

\end{enumerate}
The exclusion of \(0\) from the multiplicative group is essential.  If \(0\) had a multiplicative inverse, then from \(0\cdot a=0\) one would get contradictions such as \(1=0\).  So the phrase ``every nonzero element has an inverse'' is not a technical nuisance; it is what makes division possible.

Examples are
\[
  \mathbb Q,\qquad \mathbb R,\qquad \mathbb C,
\]
and, for a prime \(p\),
\[
  \mathbb F_p=\mathbb Z/p\mathbb Z.
\]
The requirement that \(p\) be prime is crucial.  If \(p\) is composite, zero divisors appear.  For example in \(\mathbb Z/6\mathbb Z\),
\[
  2\cdot3\equiv0\pmod6,
\]
even though neither factor is zero.

\section{Field extensions}

A field extension \(K/F\) means \(F\) is embedded in a larger field \(K\).  The point of passing to \(K\) is often to make more polynomial equations solvable.  For example,
\[
  x^2+1
\]
has no root in \(\mathbb R\), but it splits in \(\mathbb C\):
\[
  x^2+1=(x-i)(x+i).
\]

If \(K\) is finite-dimensional as a vector space over \(F\), its dimension is the degree of the extension:
\[
  [K:F]=\dim_F K.
\]
The note's intuition is good: finite extensions are tractable because linear algebra can be used.  Infinite extensions require much heavier tools.

\section{Constructing finite fields}

For a prime \(p\), \(\mathbb F_p\) is the simplest finite field.  To construct fields with \(p^n\) elements, use irreducible polynomials.  If
\[
  f(x)\in\mathbb F_p[x]
\]
is irreducible of degree \(n\), then
\[
  \mathbb F_p[x]/(f(x))
\]
is a field with \(p^n\) elements.

For example, over \(\mathbb F_2\), the polynomial
\[
  f(x)=x^3+x+1
\]
has no root in \(\mathbb F_2\), hence is irreducible of degree \(3\).  Therefore
\[
  \mathbb F_2[x]/(x^3+x+1)
\]
is a field with \(2^3=8\) elements.

In this quotient field, if \(\alpha\) denotes the class of \(x\), then
\[
  \alpha^3+\alpha+1=0,
\]
so
\[
  \alpha^3=\alpha+1
\]
in characteristic \(2\).  Every element can be uniquely written as
\[
  a+b\alpha+c\alpha^2,\qquad a,b,c\in\mathbb F_2.
\]

\section{Important properties of finite fields}

The standard facts are:
\begin{enumerate}[(i)]
  \item finite fields have order \(p^n\);
  \item for every prime power \(p^n\), there is a finite field of that order, unique up to isomorphism;
  \item the nonzero elements \(\mathbb F_{p^n}^\times\) form a cyclic group of order \(p^n-1\);
  \item the Frobenius map
  \[
    \operatorname{Frob}(a)=a^p
  \]
  is an automorphism of \(\mathbb F_{p^n}\).

\end{enumerate}

The cyclicity of the multiplicative group is surprisingly powerful.  It means every nonzero element is a power of one primitive element.  This is why finite field computations often become exponent arithmetic modulo \(p^n-1\).

\section{Polynomial factorization over finite fields}

In \(\mathbb F_q[x]\), every nonzero polynomial factors uniquely into irreducible polynomials, up to order and multiplication by nonzero constants.  This is analogous to unique prime factorization of integers.

However, a polynomial need not split into linear factors over its base finite field.  For example, a polynomial may be irreducible over \(\mathbb F_p\) but split over an extension field.  This is exactly like \(x^2+1\) over \(\mathbb R\), except that the extension fields are finite and highly structured.

\section[Example over F\_3]{Example over \(\mathbb F_3\)}

Consider
\[
  f(x)=x^4+x^2+2\in\mathbb F_3[x].
\]
First test roots:
\[
  f(0)=2\neq0,
\]
\[
  f(1)=1+1+2=4\equiv1\pmod3,
\]
\[
  f(2)=2^4+2^2+2=16+4+2=22\equiv1\pmod3.
\]
So \(f\) has no linear factor.  It might still factor as a product of two quadratics.  A systematic check over \(\mathbb F_3\) shows that no such factorization exists, so \(f\) is irreducible over \(\mathbb F_3\).

This illustrates the basic workflow: root testing only detects linear factors.  For degree \(4\), one must also check quadratic factors.

\section{Algorithms}

The note lists several factorization algorithms:
\begin{enumerate}[(i)]
  \item trial division, useful for very small degrees;
  \item Berlekamp's algorithm, which converts factorization into linear algebra over finite fields;
  \item Cantor--Zassenhaus, a randomized algorithm widely used in computer algebra systems.

\end{enumerate}
The high-level idea behind these algorithms is that finite fields make the Frobenius map computable.  Since every element of \(\mathbb F_q\) satisfies
\[
  a^q=a,
\]
the polynomial
\[
  x^q-x
\]
contains all elements of \(\mathbb F_q\) as roots.  This identity becomes a machine for detecting factors.

\section{Applications}

Finite-field factorization appears in:
\begin{enumerate}[(i)]
  \item error-correcting codes;
  \item cryptographic protocols;
  \item construction of extension fields;
  \item computational number theory;
  \item symbolic algebra systems.

\end{enumerate}
The note's final reflection is that finite fields are small enough for computation but structured enough to retain deep algebraic behavior.  That is why they are so central.


\section{Irreducibility over finite fields}

A polynomial can factor differently depending on the coefficient field.  For example, \(x^2+1\) is irreducible over \(\mathbb R\) only if one insists on real linear factors, but over \(\mathbb C\) it splits.  Over finite fields the same phenomenon becomes arithmetic.

For \(\mathbb F_p\), a quadratic polynomial
\[
  ax^2+bx+c
\]
factors iff its discriminant
\[
  \Delta=b^2-4ac
\]
is a square in \(\mathbb F_p\).  Thus irreducibility is controlled by quadratic residues.  This links polynomial factorization directly to modular arithmetic.

\section{Constructing extensions}

If \(f(x)\in F[x]\) is irreducible of degree \(n\), then
\[
  F[x]/(f(x))
\]
is a field extension of \(F\) of degree \(n\).  The class of \(x\) becomes an element \(\alpha\) satisfying \(f(\alpha)=0\).  This is the algebraic meaning of “adjoining a root.”

For finite fields, this construction produces fields with \(p^n\) elements:
\[
  \mathbb F_{p^n}\cong \mathbb F_p[x]/(f(x))
\]
for any irreducible \(f\) of degree \(n\).  Every finite field has prime-power size, and for each prime power there is a unique finite field up to isomorphism.

\section{Frobenius map}

In characteristic \(p\), the Frobenius map is
\[
  \varphi(a)=a^p.
\]
It is a field homomorphism because
\[
  (a+b)^p=a^p+b^p
\]
in characteristic \(p\).  In a finite field, Frobenius is an automorphism.  Its fixed field inside \(\mathbb F_{p^n}\) is \(\mathbb F_p\).

The polynomial
\[
  x^{p^n}-x
\]
has all elements of \(\mathbb F_{p^n}\) as roots, and it factors over \(\mathbb F_p\) as the product of all monic irreducible polynomials whose degrees divide \(n\).  This is a central organizing fact for finite-field factorization.

\section{Expanded synthesis and advanced context}

Finite fields turn algebra into arithmetic.  Roots of polynomials become elements of finite extensions, irreducibility becomes a modular phenomenon, and Frobenius gives a canonical symmetry.  This is why finite fields appear in coding theory and cryptography: their algebra is rich but completely finite and computable.


\paragraph{Expanded synthesis.}
For this chapter, the elementary material should be read as a study of what remains invariant when coordinates change. The earlier sections (`Why finite fields matter`, `Fields`, `Field extensions`, `Constructing finite fields`, `Important properties of finite fields`) compute with arrays, bases, pivots, determinants, or eigenvectors; the deeper object is a linear map together with the spaces it acts on. A matrix is a representation of that map after bases have been chosen. This is why formulas such as
\[
  A' = P^{-1}AP,\qquad \widetilde A = T^{-1}AS
\]
are not cosmetic: they distinguish an intrinsic morphism from a coordinate description.

The structural hierarchy is as follows. Kernels record what a map forgets, images record what it reaches, cokernels record obstructions to solvability, and quotients record information after an equivalence has been imposed. Determinants act on the top exterior power and therefore measure oriented volume; traces are invariant under cyclic rearrangement and become characters in representation theory; adjoints depend on an inner product and lead to self-adjoint spectral theory.

A useful way to return to the computations is to ask, for every formula, which invariant it is measuring. Row reduction measures rank and kernel; diagonalization measures decomposition into invariant subspaces; Gram--Schmidt measures orthogonal projection; positive definiteness measures ellipsoid geometry; tensor notation measures how components transform. The elementary methods are therefore the finite-dimensional laboratory in which duality, quotienting, functoriality, and spectral decomposition can be seen clearly.


\section{Expanded synthesis and advanced context}

A finite field is where group theory, ring theory, linear algebra, and number theory meet.  Addition gives a vector space over \(\mathbb F_p\); multiplication gives a cyclic group on nonzero elements; Frobenius gives a canonical symmetry; irreducible polynomials give field extensions.  Polynomial factorization over finite fields is therefore not just a computational task, but a compact model of algebra itself.



\paragraph{Expanded synthesis.}
For this chapter, the elementary material should be read as a study of what remains invariant when coordinates change. The earlier sections (`Why finite fields matter`, `Fields`, `Field extensions`, `Constructing finite fields`, `Important properties of finite fields`) compute with arrays, bases, pivots, determinants, or eigenvectors; the deeper object is a linear map together with the spaces it acts on. A matrix is a representation of that map after bases have been chosen. This is why formulas such as
\[
  A' = P^{-1}AP,\qquad \widetilde A = T^{-1}AS
\]
are not cosmetic: they distinguish an intrinsic morphism from a coordinate description.

The structural hierarchy is as follows. Kernels record what a map forgets, images record what it reaches, cokernels record obstructions to solvability, and quotients record information after an equivalence has been imposed. Determinants act on the top exterior power and therefore measure oriented volume; traces are invariant under cyclic rearrangement and become characters in representation theory; adjoints depend on an inner product and lead to self-adjoint spectral theory.

A useful way to return to the computations is to ask, for every formula, which invariant it is measuring. Row reduction measures rank and kernel; diagonalization measures decomposition into invariant subspaces; Gram--Schmidt measures orthogonal projection; positive definiteness measures ellipsoid geometry; tensor notation measures how components transform. The elementary methods are therefore the finite-dimensional laboratory in which duality, quotienting, functoriality, and spectral decomposition can be seen clearly.


\section{Elementary-to-advanced bridge: Frobenius automorphism}

In a finite field of characteristic $p$, the Frobenius map is
\[
  F:x\mapsto x^p.
\]
It is a field homomorphism because
\[
  (x+y)^p=x^p+y^p
\]
in characteristic $p$.  In a finite field it is an automorphism.

For $\mathbb F_{p^n}$, the Galois group over $\mathbb F_p$ is cyclic generated by Frobenius:
\[
  \operatorname{Gal}(\mathbb F_{p^n}/\mathbb F_p)=\langle F\rangle,\qquad F^n=1.
\]
This is the cleanest first example of a Galois group.

\section{Elementary-to-advanced bridge: irreducible-polynomial counting}

Elements of $\mathbb F_{q^n}$ are roots of
\[
  x^{q^n}-x.
\]
This polynomial factors over $\mathbb F_q$ as the product of all monic irreducible polynomials whose degrees divide $n$.

Let $N_d$ be the number of monic irreducible polynomials of degree $d$ over $\mathbb F_q$.  Then
\[
  q^n=\sum_{d\mid n} dN_d.
\]
By Möbius inversion,
\[
  N_n=\frac1n\sum_{d\mid n}\mu(d)q^{n/d}.
\]
This is a direct bridge from elementary polynomial factorization to finite-field enumeration.

\section{Elementary-to-advanced bridge: coding theory}

Finite fields are the alphabet for many error-correcting codes.  A linear code is a subspace
\[
  C\subseteq \mathbb F_q^n.
\]
The Hamming distance counts positions in which two codewords differ.  A code with large minimum distance can detect and correct errors.

Reed--Solomon codes evaluate low-degree polynomials over finite fields:
\[
  f(x)\mapsto (f(a_1),\ldots,f(a_n)).
\]
The fact that a low-degree polynomial cannot have too many roots is exactly what gives error correction.

\section{Elementary-to-advanced bridge: algebraic geometry over finite fields}

Solving polynomial equations over finite fields leads to algebraic varieties over finite fields.  Counting their points is a central research problem.  The simplest example is a curve $C$ over $\mathbb F_q$, where one studies
\[
  |C(\mathbb F_q)|.
\]
The Weil bounds relate point counts to topology-like invariants of the curve.

Thus elementary finite-field polynomial factorization is the first step toward modern arithmetic geometry.

% Second-pass deepening for waves 2-4: wave 3 A-C part 3
\section{Second-pass deepening: finite fields as vector spaces plus cyclic multiplication}

A finite field $\mathbb F_{p^n}$ is simultaneously an $n$-dimensional vector space over $\mathbb F_p$ and a field.  Its additive structure is
\[
  (\mathbb Z/p\mathbb Z)^n,
\]
while its nonzero multiplicative group is cyclic of order $p^n-1$.

This dual nature explains many finite-field constructions.  Addition is linear algebra; multiplication is encoded by reducing polynomials modulo an irreducible polynomial; Frobenius is both field automorphism and $\mathbb F_p$-linear map.

\section{Second-pass deepening: irreducible polynomials as field-building data}

Choosing an irreducible polynomial $f(x)$ of degree $n$ over $\mathbb F_p$ constructs
\[
  \mathbb F_p[x]/(f(x)).
\]
A root $\alpha$ of $f$ becomes an element of degree $n$, and every field element is represented uniquely as
\[
  a_0+a_1\alpha+\cdots+a_{n-1}\alpha^{n-1}.
\]
The original factorization exercises are therefore not only about breaking polynomials apart; irreducible polynomials are the atoms from which finite field extensions are built.
